\documentclass[12pt,a4paper]{report}

% --- PACKAGES ---
\usepackage[french]{babel}
\usepackage[T1]{fontenc}
\usepackage[utf8]{inputenc}
\usepackage{listings}
\usepackage{xcolor}
\usepackage[most]{tcolorbox}
\usepackage{geometry}

\geometry{margin=2.5cm}

% --- COULEURS ---
\definecolor{display-green}{rgb}{0.1, 0.5, 0.1}
\definecolor{code-bg}{rgb}{0.98, 0.98, 0.98}

% --- STYLE DES BOITES (Design corrigé pour Yahya) ---
\newtcolorbox{displaybox}[1]{
    enhanced,
    attach boxed title to top left={yshift=-2mm, xshift=2mm},
    colback=white,
    colframe=display-green,
    fonttitle=\bfseries,
    colbacktitle=display-green,
    title=#1,
    boxed title style={size=small, sharp corners},
    drop shadow,
    bottom=2mm
}

% --- CONFIGURATION C ---
\lstset{
    language=C,
    backgroundcolor=\color{code-bg},
    basicstyle=\ttfamily\footnotesize,
    keywordstyle=\color{blue},
    commentstyle=\color{gray},
    stringstyle=\color{purple},
    breaklines=true,
    frame=none,
    numbers=left,
    numberstyle=\tiny\color{gray}
}

\begin{document}

\chapter{Module d'Affichage et Feedback (Display)}

\section{Introduction}
Le module d'affichage regroupe tous les widgets dits de "lecture seule". Ils permettent de présenter des informations textuelles, des indicateurs de progression ou des éléments graphiques (icônes et images) à l'utilisateur sans qu'il n'ait à interagir directement avec eux.

\section{Composants de Texte et État (Header)}

\begin{displaybox}{Structure : progress\_bar\_config}
La barre de progression est indispensable pour les tâches de fond. Elle supporte un mode fractionné ($0.0$ à $1.0$) ou un mode "pulse" pour indiquer une activité sans durée définie.
\end{displaybox}

\begin{lstlisting}
typedef struct {
    double fraction;
    const char *text;
    GtkOrientation orientation;
    bool show_text;
    bool inverted;
    bool pulse_once;
    widget_style_config style;
} progress_bar_config;

#define PROGRESS_BAR_CONFIG_DEFAULT ((progress_bar_config){ \
    .fraction = 0.0, \
    .orientation = GTK_ORIENTATION_HORIZONTAL, \
    .style = WIDGET_STYLE_CONFIG_DEFAULT, \
})
\end{lstlisting}

\begin{displaybox}{Structure : image\_config}
Ce wrapper unifie deux widgets GTK différents : \texttt{GtkImage} pour les icônes système et \texttt{GtkPicture} pour les fichiers images externes, tout en gérant le respect du ratio d'aspect.
\end{displaybox}

\begin{lstlisting}
typedef struct {
    const char *icon_name;
    const char *file_path;
    int pixel_size;
    bool keep_aspect;
    widget_style_config style;
} image_config;
\end{lstlisting}

\section{Logique de Création (Source)}

\begin{displaybox}{Fonction : create\_label}
Cette fonction simplifie la création de labels. Elle gère automatiquement le retour à la ligne (\textit{wrap}) et l'élision (\textit{ellipsize}) si le texte est trop long pour l'espace alloué.
\end{displaybox}

\begin{lstlisting}
GtkWidget *create_label(const label_config *config) {
    GtkWidget *label = gtk_label_new(config && config->text ? config->text : "");
    if (config == NULL) return label;

    gtk_label_set_selectable(GTK_LABEL(label), config->selectable);
    gtk_label_set_wrap(GTK_LABEL(label), config->wrap);
    
    apply_widget_style(label, &config->style);
    return label;
}
\end{lstlisting}

\section{Gestion des Médias Visuels}



\begin{displaybox}{Fonction : create\_image}
La fonction fait un choix intelligent : si un chemin de fichier est fourni, elle utilise un \texttt{GtkPicture} (plus performant pour les photos). Sinon, elle utilise un \texttt{GtkImage} pour puiser dans le thème d'icônes du système.
\end{displaybox}

\begin{lstlisting}
GtkWidget *create_image(const image_config *config) {
    GtkWidget *widget;
    if (config == NULL) return gtk_image_new();

    if (config->file_path != NULL && config->file_path[0] != '\0') {
        GtkWidget *pic = gtk_picture_new_for_filename(config->file_path);
        gtk_picture_set_content_fit(pic, config->keep_aspect ? 
            GTK_CONTENT_FIT_CONTAIN : GTK_CONTENT_FIT_FILL);
        widget = pic;
    } else {
        widget = gtk_image_new_from_icon_name(config->icon_name);
    }

    apply_widget_style(widget, &config->style);
    return widget;
}
\end{lstlisting}

\end{document}